% Bölüm 22 — taslak plan. Tam metin yazıldığında bu dosyanın yerini alır.
\chapter{Yapay Zekâ ile Teorem İspatı}\label{bol:yapay-zeka}

\begin{bolumhedefleri}
\begin{itemize}
  \item Büyük dil modellerinin formal ispattaki rolünü
  \item Kıyaslama setlerini: miniF2F, PutnamBench
  \item Güncel sistemleri ve sınırlarını
  \item Formal doğrulamanın yapay zekâ güvenliğine katkısını
\end{itemize}
\end{bolumhedefleri}

\begin{taslakplan}
\textbf{Planlanan kesimler}
\begin{enumerate}[label=\thechapter.\arabic*, leftmargin=2.4em]
  \item Neden formal ispat? Halüsinasyona karşı doğrulama
  \item Kıyaslama setleri
  \item LeanDojo ve erişim destekli ispat
  \item AlphaProof, Seed-Prover ve sonraki sistemler
  \item Otomatik formalizasyon (autoformalization)
  \item Açık problemler
\end{enumerate}

\textbf{Planlanan Lean örnekleri}
\begin{itemize}
  \item \lc{Bir yapay zekâ ispatını Lean ile denetlemek}
\end{itemize}

\textbf{Temel kaynaklar:} \citet{zheng2022minif2f,yang2023leandojo,tsoukalas2024putnambench,alphaproof2024,seedprover2025}.
\end{taslakplan}

